%%%%%%%%%%%%%%%%%%%%%%%%%%%%%%%%%%%%%%%%%%%%%%%%%%%%%%%%%%%%%%%%%
% Presentacion 3/3: Cadenas, Archivos y Funciones Avanzadas en Python
% Serie "Python desde C" - TP4 Parte B (UNCuyo - Informatica y Programacion)
% Compilar con: pdflatex -shell-escape 03-cadenas-archivos-funciones.tex  (dos pasadas)
%%%%%%%%%%%%%%%%%%%%%%%%%%%%%%%%%%%%%%%%%%%%%%%%%%%%%%%%%%%%%%%%%
\documentclass{beamer}

% Paquetes basicos
\usepackage[utf8]{inputenc}
\usepackage[T1]{fontenc}
\usepackage[spanish]{babel}

% Paquete para el resaltado de codigo de alta calidad
\usepackage{minted}

% Tema visual de la presentacion
\usetheme{Madrid}
\usecolortheme{beaver}
\setbeamertemplate{navigation symbols}{}

% Recapitula el indice, resaltando la seccion actual, al empezar cada seccion
\AtBeginSection[]{
  \begin{frame}{Indice de Contenidos}
    \tableofcontents[currentsection]
  \end{frame}
}

% --- Informacion del Titulo ---
\title[Python desde C -- 3/3]{Cadenas, Archivos y Funciones Avanzadas en Python}
\subtitle{3/3 -- De cadenas y \texttt{string.h} a \texttt{open()} y \texttt{stdio.h}}
\author{Informática y Programación}
\date{\today}
\institute{UNCuyo -- Facultad de Ingeniería}

% --- Inicio del Documento ---
\begin{document}

% --- Diapositiva del Titulo ---
\begin{frame}
  \titlepage
\end{frame}

% --- Diapositiva del Indice ---
\begin{frame}{Índice de Contenidos}
  \tableofcontents
\end{frame}

\begin{frame}{El cierre de la serie}
  \begin{block}{Dónde estamos}
    \begin{enumerate}
      \item Fundamentos: sintaxis, tipos, control, E/S, funciones.
      \item Estructuras de datos: listas, tuplas y diccionarios.
      \item \textbf{Cadenas, archivos y funciones avanzadas} (ésta).
    \end{enumerate}
  \end{block}
  \begin{alertblock}{Para qué sirve esta última parte}
    Cierra lo que falta para el ejercicio 7 (Wordle, manipulación de
    cadenas) y resuelve directamente el ejercicio 9 del TP4: leer un
    archivo de texto y calcular estadísticas sobre su contenido.
  \end{alertblock}
\end{frame}


% ===============================================================
\section{1. Cadenas: Indexado y Slicing}
% ===============================================================

\begin{frame}[fragile]{Acceder a un carácter}
  \begin{block}{Misma idea, mismo símbolo}
    En los dos lenguajes una cadena se indexa como un arreglo: posición
    \texttt{0} es el primer carácter.
  \end{block}
  \begin{columns}[T]
    \begin{column}{.48\textwidth}
      \centering\textbf{C}
      \begin{minted}[fontsize=\small, frame=lines, linenos]{c}
char s[] = "Hola";
printf("%c\n", s[0]);
// Imprime: H
      \end{minted}
    \end{column}
    \begin{column}{.48\textwidth}
      \centering\textbf{Python}
      \begin{minted}[fontsize=\small, frame=lines, linenos]{python}
s = "Hola"
print(s[0])
# Imprime: H
      \end{minted}
    \end{column}
  \end{columns}
\end{frame}

\begin{frame}[fragile]{Obtener una subcadena: \texttt{slicing}}
  \begin{block}{Acá empiezan a distanciarse}
    Python tiene \texttt{slicing} incorporado al lenguaje:
    \texttt{s[inicio:fin]}. En C no existe — hay que armar la subcadena a
    mano, carácter por carácter.
  \end{block}
  \begin{columns}[T]
    \begin{column}{.48\textwidth}
      \centering\textbf{C}
      \begin{minted}[fontsize=\footnotesize, frame=lines, linenos]{c}
char s[] = "Programacion";
char sub[4];
for (int i = 0; i < 3; i++)
    sub[i] = s[2 + i];
sub[3] = '\0';
// sub = "ogr"
      \end{minted}
    \end{column}
    \begin{column}{.48\textwidth}
      \centering\textbf{Python}
      \begin{minted}[fontsize=\footnotesize, frame=lines, linenos]{python}
s = "Programacion"
sub = s[2:5]
# sub = "ogr"
      \end{minted}
    \end{column}
  \end{columns}
\end{frame}

\begin{frame}[fragile]{Bonus: invertir una cadena}
  \begin{block}{\texttt{s[::-1]}: paso negativo}
    El tercer valor del slice es el \textbf{paso}. Con paso \texttt{-1},
    Python recorre la cadena de atrás para adelante — útil para, por
    ejemplo, verificar palíndromos (ya lo hicieron en C, en el TP3, con
    recursión).
  \end{block}
  \begin{columns}[T]
    \begin{column}{.48\textwidth}
      \centering\textbf{C}
      \begin{minted}[fontsize=\footnotesize, frame=lines, linenos]{c}
int n = strlen(s);
for (int i = n - 1; i >= 0; i--)
    putchar(s[i]);
// recorrido manual
      \end{minted}
    \end{column}
    \begin{column}{.48\textwidth}
      \centering\textbf{Python}
      \begin{minted}[fontsize=\footnotesize, frame=lines, linenos]{python}
s = "Hola"
print(s[::-1])
# Imprime: aloH
      \end{minted}
    \end{column}
  \end{columns}
\end{frame}

% ===============================================================
\section{2. Cadenas Inmutables}
% ===============================================================

\begin{frame}{Las cadenas en Python no se pueden modificar}
  \begin{alertblock}{Contraintuitivo viniendo de C}
    En C, \texttt{s} es un arreglo de \texttt{char}: cada posición se
    puede reescribir libremente (\texttt{s[0] = 'X';}). En Python, un
    \texttt{str} es \textbf{inmutable} — una vez creado, ningún carácter
    se puede cambiar en el lugar. Cualquier operación que ``modifica'' una
    cadena en realidad \textbf{crea una cadena nueva}.
  \end{alertblock}
  \begin{block}{Por qué importa}
    Es la razón por la que \texttt{s = s + "algo"} o \texttt{s[2:5]} no
    tocan la cadena original: siempre devuelven una cadena distinta.
  \end{block}
\end{frame}

\begin{frame}[fragile]{El error típico y cómo evitarlo}
  \begin{columns}[T]
    \begin{column}{.48\textwidth}
      \centering\textbf{C (válido)}
      \begin{minted}[fontsize=\footnotesize, frame=lines, linenos]{c}
char s[] = "Hola";
s[0] = 'h';
printf("%s\n", s);
// Imprime: hola
      \end{minted}
    \end{column}
    \begin{column}{.48\textwidth}
      \centering\textbf{Python (error)}
      \begin{minted}[fontsize=\footnotesize, frame=lines, linenos]{python}
s = "Hola"
s[0] = 'h'
# TypeError: 'str' object
# does not support item
# assignment
      \end{minted}
    \end{column}
  \end{columns}
  \vspace{0.4em}
  \begin{minted}[fontsize=\footnotesize, frame=lines, linenos]{python}
s = "Hola"
s = "h" + s[1:]   # se arma una cadena nueva
print(s)           # Imprime: hola
  \end{minted}
\end{frame}

% ===============================================================
\section{3. Métodos de Cadena}
% ===============================================================

\begin{frame}{Lo que en C es una función suelta, en Python es un método}
  \begin{block}{\texttt{cadena.metodo()} en vez de \texttt{funcion(cadena)}}
    \texttt{\textless string.h\textgreater} da funciones que reciben la cadena como parámetro;
    Python las expone como métodos del propio \texttt{str}.
  \end{block}
  \begin{center}
  \renewcommand{\arraystretch}{1.25}
  \begin{tabular}{|l|l|}
    \hline
    \textbf{Python} & \textbf{Equivalente aproximado en C} \\
    \hline \hline
    \texttt{s.upper()} / \texttt{s.lower()} & bucle manual con \texttt{toupper}/\texttt{tolower} \\
    \hline
    \texttt{s.strip()} & recorrer y descartar espacios en los bordes \\
    \hline
    \texttt{s.split(sep)} & \texttt{strtok()} (con estado global, más delicado) \\
    \hline
    \texttt{sep.join(lista)} & concatenar a mano con \texttt{strcat()} \\
    \hline
    \texttt{s.count(sub)} & bucle manual buscando ocurrencias \\
    \hline
    \texttt{s.replace(a, b)} & sin equivalente directo, se arma a mano \\
    \hline
  \end{tabular}
  \end{center}
\end{frame}

\begin{frame}[fragile]{\texttt{split()} y \texttt{join()}: separar y armar}
  \begin{block}{Los inversos el uno del otro}
    \texttt{split()} separa una cadena en una lista; \texttt{join()} hace
    lo opuesto: une una lista en una única cadena, con un separador.
  \end{block}
  \begin{minted}[fontsize=\small, frame=lines, linenos]{python}
frase = "el perro come carne"
palabras = frase.split()          # ['el', 'perro', 'come', 'carne']

reconstruida = "-".join(palabras) # "el-perro-come-carne"
  \end{minted}
  {\footnotesize \textit{En C, la función equivalente es \texttt{strtok()}
  — funciona, pero modifica la cadena original y guarda estado entre
  llamadas, lo que la hace más propensa a errores que \texttt{split()}.}}
\end{frame}

% ===============================================================
\section{4. Manejo de Archivos}
% ===============================================================

\begin{frame}{Abrir y cerrar: \texttt{open()} + \texttt{with} vs. \texttt{fopen}/\texttt{fclose}}
  \begin{alertblock}{El mismo riesgo que con \texttt{malloc}/\texttt{free}}
    En C, todo \texttt{fopen()} exitoso necesita su \texttt{fclose()} —
    olvidarlo dejar el archivo abierto es un descuido parecido a olvidar
    un \texttt{free()} en el TP3 (Parte 2).
  \end{alertblock}
  \begin{block}{El bloque \texttt{with} de Python}
    \texttt{with open(...) as f:} cierra el archivo automáticamente al
    salir del bloque, \textbf{incluso si ocurre un error adentro}. No hay
    equivalente automático en C: el \texttt{fclose()} hay que escribirlo
    siempre a mano.
  \end{block}
\end{frame}

\begin{frame}[fragile]{Abrir, leer todo el contenido y cerrar}
  \begin{columns}[T]
    \begin{column}{.48\textwidth}
      \centering\textbf{C}
      \begin{minted}[fontsize=\footnotesize, frame=lines, linenos]{c}
FILE *f = fopen("a.txt", "r");
if (f == NULL) {
    return 1; // error
}
char buf[256];
fread(buf, 1, 255, f);
fclose(f); // obligatorio
      \end{minted}
    \end{column}
    \begin{column}{.48\textwidth}
      \centering\textbf{Python}
      \begin{minted}[fontsize=\footnotesize, frame=lines, linenos]{python}
with open("a.txt") as f:
    contenido = f.read()

# el archivo ya esta
# cerrado aca afuera
      \end{minted}
    \end{column}
  \end{columns}
\end{frame}

\begin{frame}[fragile]{Leer línea por línea}
  \begin{block}{Iterar el archivo directamente}
    Un archivo abierto con \texttt{open()} es iterable: cada vuelta del
    \texttt{for} entrega la siguiente línea, sin manejar buffers a mano.
  \end{block}
  \begin{columns}[T]
    \begin{column}{.48\textwidth}
      \centering\textbf{C}
      \begin{minted}[fontsize=\footnotesize, frame=lines, linenos]{c}
char linea[256];
FILE *f = fopen("a.txt", "r");
while (fgets(linea, 256, f)) {
    printf("%s", linea);
}
fclose(f);
      \end{minted}
    \end{column}
    \begin{column}{.48\textwidth}
      \centering\textbf{Python}
      \begin{minted}[fontsize=\footnotesize, frame=lines, linenos]{python}
with open("a.txt") as f:
    for linea in f:
        print(linea, end="")
      \end{minted}
    \end{column}
  \end{columns}
\end{frame}

% ===============================================================
\section{5. Funciones Avanzadas}
% ===============================================================

\begin{frame}[fragile]{Parámetros con valor por defecto}
  \begin{block}{Algo que C no tiene}
    Python permite que un parámetro tenga un valor si quien llama no lo
    especifica. En C no existe esto: todos los argumentos se pasan
    siempre, o se define una función distinta para cada caso.
  \end{block}
  \begin{minted}[fontsize=\small, frame=lines, linenos]{python}
def calcular_tarifa(horas, tarifa_por_hora=1000):
    return horas * tarifa_por_hora

calcular_tarifa(3)         # usa tarifa_por_hora=1000 -> 3000
calcular_tarifa(3, 2000)   # la pisa -> 6000
  \end{minted}
  {\footnotesize \textit{Repaso rápido: el retorno múltiple con tuplas
  (\texttt{return a, b}) ya lo vimos en la presentación 2 — es la otra
  herramienta clave para funciones más expresivas que en C.}}
\end{frame}

% ===============================================================
\section{6. Caso Aplicado y Resumen}
% ===============================================================

\begin{frame}[fragile]{Caso aplicado: palabra más repetida (TP4, ej. 9)}
  \begin{block}{Contar frecuencia de palabras con un diccionario}
    Recorrer las palabras y, por cada una, sumar 1 a su contador en un
    diccionario — \texttt{dict.get(clave, 0)} devuelve \texttt{0} si la
    clave todavía no existe, evitando un \texttt{if} extra.
  \end{block}
  \begin{minted}[fontsize=\footnotesize, frame=lines, linenos]{python}
frecuencia = {}
for palabra in palabras:
    if len(palabra) > 6:
        frecuencia[palabra] = frecuencia.get(palabra, 0) + 1

# la(s) palabra(s) con mayor frecuencia:
maximo = max(frecuencia.values())
mas_repetidas = [p for p, c in frecuencia.items() if c == maximo]
  \end{minted}
\end{frame}

\begin{frame}{Tabla Comparativa Final}
  \begin{center}
  \renewcommand{\arraystretch}{1.3}
  \begin{tabular}{|l|p{4cm}|p{4cm}|}
    \hline
    \textbf{Característica} & \textbf{Lenguaje C} & \textbf{Lenguaje Python} \\
    \hline \hline
    Cadenas & Arreglo de \texttt{char} mutable, con \texttt{'\textbackslash0'} & \texttt{str} inmutable, con métodos propios \\
    \hline
    Archivos & \texttt{fopen}/\texttt{fclose} manual & \texttt{with open(...)} con cierre automático \\
    \hline
    Manejo de recursos & Responsabilidad del programador & Garbage collector + \texttt{with} \\
    \hline
  \end{tabular}
  \end{center}
\end{frame}

\begin{frame}{Con esto ya tenés todo el TP4}
  \begin{block}{Parte A}
    Análisis, entradas/salidas y diseño de funciones — ya lo resolvieron
    en C y pseudocódigo.
  \end{block}
  \begin{block}{Parte B}
    Variables, control de flujo, funciones (presentación 1); listas,
    tuplas y diccionarios (presentación 2); cadenas, archivos y funciones
    avanzadas (ésta) — todo lo necesario para los ejercicios 5 a 9.
  \end{block}
\end{frame}

\begin{frame}
  \vfill
  \centering
  \begin{beamercolorbox}[sep=8pt,center,shadow=true,rounded=true]{title}
    \Huge \bfseries ¡A resolver el TP4!
  \end{beamercolorbox}
  \vfill
\end{frame}

\end{document}
